\section{Score rounding and the integrality gap}
\label{s-integrality-gap}

Fix an integer maximizer
\[
    N_m^\star=\prod_{i=1}^{m}p_i^{a_i},
\]
and put \(Y=\log(1+T^\star)\). We first show that its exponents are
obtained by rounding the relaxed exponents one coordinate at a time.

\begin{proposition}[One-step proximity]
\label{prop-floor-ceiling}
For every sufficiently large \(m\),
\[
    a_i
    \in
    \left\{
        \left\lfloor\frac{Y}{\log p_i}-1\right\rfloor,\,
        \left\lceil\frac{Y}{\log p_i}-1\right\rceil
    \right\}
\]
when \(p_i<e^{Y/2}\). For \(p_i\geq e^{Y/2}\), we have \(a_i=1\).
\end{proposition}

\begin{proof}
Fix a prime below \(e^{Y/2}\), and put
\[
    \ell=\left\lfloor\frac{Y}{\log p}-1\right\rfloor.
\]
The relaxed first-order equation is
\[
    \frac1{p^{Y/\log p}-1}=\frac1{T^\star}.
\]
If \(\ell\geq2\), the increment to exponent \(\ell\) lies below the
relaxed coordinate. Its score is greater than \(1/T^\star\), and hence greater than
\(\eps_m\) by Proposition~\ref{prop-integer-threshold-scale}. It is
selected. When \(\ell=1\), the support constraint already gives
\(a_i\geq\ell\).

We next rule out the increment to exponent \(\ell+2\). Write
\[
    u=\frac{Y}{\log p}-1+v
\]
on the corresponding unit interval. Since \(p^{Y/\log p}=1+T^\star\),
\[
    \frac{T^\star}{p^{u+1}-1}
    \leq p^{-v}.
\]
The interval begins at some \(v\in(0,1]\), so
\[
    T^\star F(p,\ell+2)
    \leq
    \int_0^1p^{-v}\,dv
    =
    \frac{1-p^{-1}}{\log p}
    \leq
    \frac1{2\log2}
    <1.
\]
Since \(T^\star\eps_m\) tends to one, this increment and all later ones
are unselected. Thus \(\ell\leq a_i\leq\ell+1\).

If \(Y/\log p-1\) is an integer, the same estimate applied to the next
increment shows that it is unselected. This proves the stated
floor--ceiling alternative. If \(p_i\geq e^{Y/2}\), apply the same
estimate to \(F(p_i,2)\). The corresponding unit interval begins at a
nonnegative value of \(v\), so this increment is unselected and
\(a_i=1\).
\end{proof}

\begin{theorem}[Integrality-gap asymptotic]
\label{thm-integrality-gap}
The integrality gap satisfies
\[
    \Igap(m)
    =
    \frac{2\sqrt2+o(1)}
         {\sqrt{p_m}\log p_m}.
\]
\end{theorem}

To compare the optimizers coordinatewise, define a score-rounded integer
exponent at every log-prime \(x\). For
\(Y>2\log2\) and \(\log2\leq x\leq Y/2\), put
\begin{equation}
\label{e-rounded-exponent}
    a_Y(x)
    =
    \begin{cases}
    \lfloor Y/x\rfloor,
        &F(e^x,\lfloor Y/x\rfloor)>(e^Y-1)^{-1},\\
    \lfloor Y/x\rfloor-1,
        &F(e^x,\lfloor Y/x\rfloor)\leq(e^Y-1)^{-1}.
    \end{cases}
\end{equation}
This is one of the two integers adjacent to \(Y/x-1\). It rounds upward
exactly when the intervening chord score exceeds the relaxed multiplier.
Define its local loss by
\begin{equation}
\label{e-kappa}
    \kappa(x,Y)
    =
    \log\frac{1-e^{-Y}}
              {1-e^{-x(a_Y(x)+1)}}
    +
    \frac{x(a_Y(x)+1)-Y}{e^Y-1},
\end{equation}
and define
\begin{equation}
\label{e-rounding-integral}
    J(Y)
    =
    \int_{\log2}^{Y/2}
    \kappa(x,Y)\frac{e^x}{x}\,dx.
\end{equation}
Thus \(\kappa\) is the loss in the local divisor factor after the
first-order log-size change has been removed.

\begin{proof}[Proof of Theorem~\ref{thm-integrality-gap}]
\medskip\noindent\emph{Rounding reduction.}
For \(p_i<e^{Y/2}\), let \(b_i=a_Y(\log p_i)\). Subtracting the
two objective values and linearizing the size penalty at \(S^\star\)
gives
\[
\begin{aligned}
    \Igap(m)
    &=
    \sum_{p_i<e^{Y/2}}
    \left[
        \log\frac{p_i-p_i^{-r_i^\star}}
                  {p_i-p_i^{-a_i}}
        -\frac{(r_i^\star-a_i)\log p_i}{T^\star}
    \right]\\
    &\quad
    +O\left(
        \frac{\log T^\star}{(T^\star)^2}
        (L_m-S^\star)^2
    \right).
\end{aligned}
\]
Proposition~\ref{prop-floor-ceiling} gives
\[
    |L_m-S^\star|
    \leq
    \theta(e^{Y/2})
    =
    O(\sqrt{T^\star}),
\]
so the displayed error is lower order on the required scale.

Both \(a_i\) and \(b_i\) are adjacent to \(r_i^\star\). If they differ,
the intervening score lies between \(\eps_m\) and \(1/T^\star\). Hence
replacing \(a_i\) by \(b_i\) in the sum changes it by at most
\[
    \theta(e^{Y/2})
    \left(\frac1{T^\star}-\eps_m\right)
    =
    O\left(\frac{\sqrt{\log T^\star}}{T^\star}\right),
\]
where we used \eqref{e-threshold-separation}. This is also lower order.
The bracket with \(b_i\) is exactly
\(\kappa(\log p_i,Y)\), by \eqref{e-kappa}.
We have proved
\[
    \sum_{p<e^{Y/2}}\kappa(\log p,Y)
    =
    \Igap(m)
    +
    o\left(\frac1{\sqrt{T^\star\log T^\star}}\right).
\]

\medskip\noindent\emph{Prime-sum approximation.}
Put \(q=e^{-Y}\) and \(\lambda=(e^Y-1)^{-1}\). On an interval where
\(k=\lfloor Y/x\rfloor\) is constant, put \(s=Y-kx\). Formula
\eqref{e-kappa} becomes
\begin{equation}
\label{e-kappa-branches}
    \kappa(x,Y)
    =
    \begin{cases}
    \log(1-q)-\log(1-qe^s)-s\lambda,
        &F(e^x,k)\leq\lambda,\\[1.5mm]
    \log(1-q)-\log(1-qe^{s-x})+(x-s)\lambda,
        &F(e^x,k)>\lambda.
    \end{cases}
\end{equation}
The difference between the branches is
\[
    x\bigl(\lambda-F(e^x,k)\bigr),
\]
so they agree at a switch.

On the first branch,
\[
\begin{aligned}
    xF(e^x,k)
    &=
    \log(1-qe^{s-x})-\log(1-qe^s)\\
    &\geq
    qe^s(1-e^{-x})
    \geq
    \tfrac12qe^s.
\end{aligned}
\]
The branch condition therefore gives \(qe^s\ll xq\). On the second
branch, \(qe^{s-x}\leq q\). It follows from
\eqref{e-kappa-branches} that
\[
    |\kappa(x,Y)|\ll Y e^{-Y}.
\]
Within a shell, \(s'=-k\), and differentiation gives
\[
    |\partial_x\kappa(x,Y)|\ll Y e^{-Y}
\]
on every smooth piece. There are \(O(Y)\) shells. The score
\[
    F(e^x,k)=\int_k^{k+1}\frac{du}{e^{xu}-1}
\]
is \eqref{e-score-integral} after shifting the integration variable. It
decreases strictly with \(x\), so each shell has at most one branch
switch. The smooth pieces have total length \(O(Y)\), and the
shell-boundary jumps are \(O(Ye^{-Y})\). Consequently, the total
variation satisfies
\[
    \operatorname{TV}_{[\log2,Y/2]}
    \kappa(\,\cdot\,,Y)
    \ll
    Y^2e^{-Y}.
\]
At \(x=Y/2\), the score-rounded exponent is one and
\(\kappa(Y/2,Y)=0\).

The de la Vall\'ee Poussin estimate \cite{delavalleepoussin1899} gives
\[
    \pi(e^x)-\operatorname{li}(e^x)
    =
    O\left(
        \frac{e^x}{x}e^{-c\sqrt x}
    \right).
\]
Stieltjes integration by parts, using the variation bound and the
vanishing upper endpoint, now gives
\[
\begin{aligned}
    &\sum_{p<e^{Y/2}}\kappa(\log p,Y)
    -
    \int_{\log2}^{Y/2}\kappa(x,Y)\frac{e^x}{x}\,dx\\
    &\qquad
    =
    O\left(
        Ye^{-Y/2}e^{-c'\sqrt Y}
    \right)
    +O(Ye^{-Y}).
\end{aligned}
\]
This is lower order than the required scale. Shell-boundary conventions
change the expression by at most \(O(Y^2e^{-Y})\).
It follows that
\[
    \sum_{p<e^{Y/2}}\kappa(\log p,Y)
    =
    J(Y)
    +
    o\left(\frac1{\sqrt{(e^Y-1)Y}}\right).
\]

\medskip\noindent\emph{Second-exponent range.}
The branch formulas give the sharper bound
\[
    |\kappa(x,Y)|\ll q(1+x).
\]
Therefore
\[
    \int_{\log2}^{Y/3}
    |\kappa(x,Y)|\frac{e^x}{x}\,dx
    =
    O(e^{-2Y/3}),
\]
which is lower order. Only \(Y/3<x\leq Y/2\), where
\(\lfloor Y/x\rfloor=2\), contributes at first order.

For all sufficiently large \(Y\), let \(x_\star\) be the solution in
\((Y/3,Y/2)\) of
\begin{equation}
\label{e-rounding-switch}
    F(e^{x_\star},2)=\frac1{e^Y-1}.
\end{equation}
On this shell put
\[
    u=Y-2x,
    \qquad
    x=\frac{Y-u}{2}.
\]
The score in \eqref{e-rounding-switch} increases with \(u\). It is below
\((e^Y-1)^{-1}\) at \(u=0\) and above it near \(u=Y/3\), so the switch
exists and is unique. Put \(u_\star=Y-2x_\star\). Uniformly on the shell,
\[
\begin{aligned}
    xF(e^x,2)
    &=
    \log(1-e^{-3x})-\log(1-e^{-2x})\\
    &=
    e^{-Y+u}
    \bigl(1+O(e^{-x}+e^{-Y+u})\bigr).
\end{aligned}
\]
The switch equation gives
\[
    e^{u_\star}
    =
    x_\star\bigl(1+O(e^{-Y/3})\bigr).
\]
Together with \(2x_\star+u_\star=Y\), this proves
\[
    2x_\star+\log x_\star=Y+o(1).
\]

On the two sides of the switch, \eqref{e-kappa-branches} gives
\[
    \kappa(x,Y)
    =
    \begin{cases}
    q(e^u-1-u)+O(q^2e^{2u}),
        &0\leq u\leq u_\star,\\[1mm]
    q(x-u-1+e^{-(x-u)})+O(Yq^2),
        &u_\star\leq u<Y/3.
\end{cases}
\]
Since \(dx=-du/2\) and \(e^x=e^{Y/2}e^{-u/2}\), the contribution from
each branch is \(qe^{Y/2}/2\) times its corresponding integral below.
For the first branch,
\[
    \int_0^{u_\star}
    \frac{e^{-u/2}}x(e^u-1-u)\,du
    =
    \frac{2+o(1)}{\sqrt{x_\star}}.
\]
Indeed, the terms \(1+u\) contribute \(O(x_\star^{-1})\). In the
remaining term, set \(t=u_\star-u\), use
\(x=x_\star+t/2\) and \(e^{u_\star}\sim x_\star\), and apply dominated
convergence.

For the second branch, set \(t=u-u_\star\). Then
\[
    x=x_\star-\frac t2,
    \qquad
    x-u=x_\star-u_\star-\frac{3t}{2}.
\]
Since
\[
    0\leq x-u-1+e^{-(x-u)}\leq x-u\leq x,
\]
the second integral equals
\[
    e^{-u_\star/2}
    \int_0^{Y/3-u_\star}
    e^{-t/2}
    \frac{x-u-1+e^{-(x-u)}}x\,dt.
\]
The ratio is bounded by one and tends to one for every fixed \(t\).
Since \(e^{-u_\star/2}\sim x_\star^{-1/2}\), dominated convergence gives
\[
    \int_{u_\star}^{Y/3}
    \frac{e^{-u/2}}x
    \bigl(x-u-1+e^{-(x-u)}\bigr)\,du
    =
    \frac{2+o(1)}{\sqrt{x_\star}}.
\]
After multiplication by the Jacobian factor \(e^{Y/2}/2\), the error
\(O(q^2e^{2u})\) on the first branch contributes
\[
    O(e^{-3Y/2}x_\star^{1/2}),
\]
while the error \(O(Yq^2)\) on the second branch contributes
\[
    O(e^{-3Y/2}x_\star^{-1/2}).
\]
Both are \(o(e^{-Y/2}x_\star^{-1/2})\), the scale of the main term.
Combining the two branches yields
\[
    J(Y)
    =
    e^{-Y/2}
    \left(
        \frac2{\sqrt{x_\star}}
        +o(x_\star^{-1/2})
    \right).
\]
Finally, \(x_\star\sim Y/2\) and
\[
    e^{-Y/2}
    =
    \frac{\sqrt{1-e^{-Y}}}{\sqrt{e^Y-1}}.
\]
Consequently,
\[
    J(Y)
    =
    \frac1{\sqrt{e^Y-1}}
    \left(
        \frac2{\sqrt{x_\star}}
        +o(x_\star^{-1/2})
    \right)
    =
    \frac{2\sqrt2+o(1)}
         {\sqrt{(e^Y-1)Y}}.
\]
The two sides of the switch contribute equally at first order. Since
\(e^Y-1=T^\star\), the rounding reduction and prime-sum approximation
give
\[
    \Igap(m)
    =
    \frac{2\sqrt2+o(1)}
         {\sqrt{T^\star Y}}.
\]
Since \(Y\sim\log T^\star\) and
\(T^\star\sim p_m\log p_m\), the denominator is asymptotic to
\(\sqrt{p_m}\log p_m\).
\end{proof}

\subsection{Relation to CA shells and the Axler--Nicolas comparison}
\label{s-score-connection}

The fixed-support optimizer is a constrained CA point, since its
optional increments are selected by the same score \(F\). It is also a
size-constrained record on its prescribed support. The classical CA shells
therefore give a second description of the integer optimizer, although
they do not provide the relaxed response or the rounding loss used in the
proof above.

Define the integer shell thresholds by
\[
    F(\xi,1)=\eps_m,
    \qquad
    F(\xi_k,k)=\eps_m.
\]
Nicolas \cite[Proposition~4.7]{nicolas-2022-sigma} proves
\[
    \xi^{1/k}\leq\xi_k\leq(k\xi)^{1/k},
    \qquad
    \xi_2\sim\sqrt{2\xi}.
\]
Since
\[
    F(\xi,1)\sim\frac1{\xi\log\xi},
\]
Proposition~\ref{prop-integer-threshold-scale} gives
\(\xi\sim L_m\sim p_m\). The exact threshold rule is
\[
    L_m-\theta(p_m)
    =
    \sum_{k\geq2}\theta\bigl(\min\{p_m,\xi_k\}\bigr).
\]
Nicolas's bounds make every support cap on the right inactive for all
sufficiently large \(m\). They then give
\[
    L_m-\theta(p_m)
    =
    \sum_{k\geq2}\theta(\xi_k)
    =
    \theta(\xi_2)+O(\xi^{1/3})
    =
    \sqrt{2p_m}+o(\sqrt{p_m}).
\]
The middle estimate follows from Nicolas's bounds for \(\xi_k\), the
prime number theorem, and the fact that only \(O(\log\xi)\) shells are
nonempty.

The relaxed switch is nearby but not identical. It is defined by
\[
    F(e^{x_\star},2)=\frac1{T^\star},
\]
whereas \(\xi_2\) uses \(\eps_m\). Equation
\eqref{e-threshold-separation} shows that
\[
    e^{x_\star}\sim\xi_2\sim\sqrt{2p_m}.
\]
Thus the standard CA second shell locates the score-rounding
transition, while the prime-sum integral resolves the loss on its two
sides.

The exact integer factorization, relative to the Mertens term used in
the next section, has the correction
\[
    \sum_{i=1}^m\log(1-p_i^{-(a_i+1)})
    -\bigl(\log\log L_m-\log\log\theta(p_m)\bigr).
\]
For \(\xi_2<p\leq p_m\), the exponent of \(p\) is one. The prime-square
tail estimate \cite[equations~(2.10)--(2.11)]{nicolas-2022-sigma} gives
\[
    \sum_{\xi_2<p\leq p_m}\log(1-p^{-2})
    =
    -\frac1{\xi_2\log\xi_2}
    +o\left(\frac1{\sqrt{p_m}\log p_m}\right)
    =
    -\frac{\sqrt2+o(1)}{\sqrt{p_m}\log p_m}.
\]
The remaining Euler factors are lower order. Primes with exponent two
lie above \(\xi_3\), so their contribution is
\[
    O\left(\sum_{p>\xi_3}p^{-3}\right)
    =
    O(\xi^{-2/3}).
\]
If \(a_i\geq3\), the next increment is unselected. Hence
\(p_i>\xi_{a_i+1}\), and Nicolas's lower bound gives
\(p_i^{a_i+1}>\xi\). There are \(O(\xi_3)=O(\xi^{1/3})\) such primes, so
their total contribution is also \(O(\xi^{-2/3})\). Consequently,
\[
    \sum_{i=1}^m\log(1-p_i^{-(a_i+1)})
    =
    -\frac{\sqrt2+o(1)}{\sqrt{p_m}\log p_m}.
\]
Taylor expansion and the shell-size formula similarly give
\[
    \log\log L_m-\log\log\theta(p_m)
    =
    \frac{\sqrt2+o(1)}{\sqrt{p_m}\log p_m}.
\]
The CA decomposition partitions the correction into an Euler-factor tail
and a log-size change, whereas the rounding integral partitions it at the
downward and upward rounding switch.

Nicolas's expansion of the Axler--Nicolas size-constrained comparison
\cite[equations~(5.3)--(5.5)]{nicolas-2025} contains the same formal pair
\[
    \frac1{\xi_2\log\xi_2},
    \qquad
    \frac{\xi_2}{\xi\log\xi}.
\]
The first is the same leading Euler-factor tail. Their second term measures
additional support-edge primes under a common size constraint. Ours measures
the Robin penalty applied to second-shell log-size. The common score and
shell equation explain the same \(2\sqrt2\) coefficient. The different
feasible sets and the different second terms prevent either theorem from
being a corollary of the other.

Since the Axler--Nicolas ratio tends to one, its relative gap and its
logarithm have the same leading term. The normalizations also agree after
the natural endpoint change. If \(P\)
is the largest prime in the primorial below \(X\), then
\(\log X\sim\theta(P)\sim P\). Hence the Axler--Nicolas scale satisfies
\[
    \frac1{\sqrt{\log X}\log\log X}
    \sim
    \frac1{\sqrt P\log P},
\]
which is the fixed-slice scale when \(P=p_m\).

There is a related relaxation interpretation. For every prime \(p\),
\[
    \frac{\sigma(p^a)}{p^a}
    \longrightarrow
    \frac1{1-p^{-1}}
    =
    \frac{p^a}{\varphi(p^a)}
\]
as \(a\) tends to infinity. Their ratio can therefore be viewed as
comparing the infinite-exponent local envelope with finite integer
exponents while also optimizing the support under a size constraint. The
present problem compares finite real and integer exponents on a fixed
support. In both
cases the first nontrivial marginal decision is the transition from
exponent one to exponent two. This common decision, rather than a formal
reduction between the problems, accounts for the matching constant.
